\section{Supporting estimator and trial checks}
\label{app:matrix}\label{sec:matrix}

\paragraph{Reference dependence.}
With nonempty stratum arms, saturated G-computation \citep{robins1986}, inverse-probability weighting
(IPW) using empirical stratum propensities, non-cross-fitted saturated
augmented IPW (AIPW) and
matching using every control all reproduce
$T(D)=\sum_g(n_g/n)(\bar Y_{1g}-\bar Y_{0g})$. The IPW weights recover each
stratum mean; the AIPW augmentation sums to zero within stratum. Ordinary least squares (OLS) with
additive stratum indicators instead weights contrasts by normalized
$n_g\hat p_g(1-\hat p_g)$ \citep{angrist1998,frisch1933}. These known identities
show why the reference determines what an exact recovery check establishes.
Exact 1:1 allocation within strata makes OLS, standardized and unadjusted
contrasts coincide; balance only in expectation does not guarantee this identity.

\paragraph{Learned-generator check.}\label{app:learned}
Four Gaussian mixtures, a multilayer-perceptron conditional Gaussian and two
parametric controls generate synthetic cohorts after fitting to simulated
three-stratum training data. Across 42 generator--size settings, saturated
G-computation has zero residual against $T(D)$ on every evaluable dataset;
this follows from the identity above. Nonsaturated fits can depart, with a
largest MLP residual of 3.3 outcome units under mixture generators.
The aggregate archive retains all scheduled/evaluable counts and exclusions.
Generator complexity alone does not make an identity-based check informative.
